% Standalone problem document, generated from the adjacent README.md.
% Compile with XeLaTeX. Mathematical content is maintained in Markdown.
\documentclass[11pt,a4paper]{article}
\usepackage[fontsize=10.5pt]{scrextend}
\usepackage[margin=22mm,headheight=15pt,headsep=9mm,footskip=11mm]{geometry}
\usepackage{fontspec}
\setmainfont{texgyrepagella}[Extension=.otf,UprightFont=*-regular,BoldFont=*-bold,ItalicFont=*-italic,BoldItalicFont=*-bolditalic]
\setsansfont{texgyreheros}[Extension=.otf,UprightFont=*-regular,BoldFont=*-bold,ItalicFont=*-italic,BoldItalicFont=*-bolditalic]
\setmonofont{texgyrecursor}[Extension=.otf,UprightFont=*-regular,BoldFont=*-bold,ItalicFont=*-italic,BoldItalicFont=*-bolditalic,Scale=MatchLowercase]
\usepackage{amsmath,amssymb}
\usepackage{unicode-math}
\setmathfont{texgyrepagella-math.otf}
\usepackage{xcolor,microtype,enumitem,needspace,fancyhdr}
\usepackage{bookmark}
\definecolor{ink}{HTML}{17344B}
\definecolor{accent}{HTML}{197D80}
\definecolor{muted}{HTML}{596773}
\hypersetup{colorlinks=true,linkcolor=accent,urlcolor=accent,pdftitle={AV-01 - Recognizing
the maximum finite number of
solutions},pdfauthor={Open Problems in Numerical Linear Algebra}}
\urlstyle{same}
\setlength{\parindent}{0pt}
\setlength{\parskip}{5pt plus 1pt minus 1pt}
\linespread{1.04}
\setlength{\emergencystretch}{3em}
\widowpenalty=10000
\clubpenalty=10000
\displaywidowpenalty=10000
\allowdisplaybreaks[1]
\setlist{nosep,leftmargin=1.5em}
\providecommand{\tightlist}{\setlength{\itemsep}{0pt}\setlength{\parskip}{0pt}}
\providecommand{\pandocbounded}[1]{#1}
\pagestyle{fancy}
\fancyhf{}
\fancyhead[L]{\sffamily\footnotesize\color{muted}OPEN PROBLEMS IN NUMERICAL LINEAR ALGEBRA}
\fancyhead[R]{\sffamily\bfseries\footnotesize\color{accent}AV-01}
\fancyfoot[L]{\parbox[b]{0.9\textwidth}{\sffamily\fontsize{8}{10}\selectfont\color{muted}Editorial ratings; a bounded search cannot certify that a problem remains unsolved.\\Literature check: 2026-09-11}}
\fancyfoot[R]{\sffamily\footnotesize\color{muted}\thepage}
\renewcommand{\headrulewidth}{0.3pt}
\renewcommand{\headrule}{\hbox to\headwidth{\color{accent}\leaders\hrule height \headrulewidth\hfill}}
\makeatletter
\renewcommand{\section}{\Needspace{5\baselineskip}\@startsection{section}{1}{0pt}{12pt plus 2pt minus 2pt}{4pt}{\sffamily\bfseries\large\color{ink}}}
\renewcommand{\subsection}{\Needspace{4\baselineskip}\@startsection{subsection}{2}{0pt}{9pt plus 1pt minus 1pt}{3pt}{\sffamily\bfseries\normalsize\color{ink}}}
\makeatother
\setcounter{secnumdepth}{0}
\begin{document}
{\sffamily\small\color{muted}Intervals and absolute value equations\par}
\vspace{3pt}
{\raggedright\hyphenpenalty=10000\exhyphenpenalty=10000\sffamily\bfseries\fontsize{21}{24}\selectfont\color{ink}Recognizing
the maximum finite number of solutions\par}
\vspace{7pt}
{\sffamily\small\textbf{Difficulty:} challenging\quad\textbf{Importance:} interesting
to specialist\par}
{\sffamily\small\bfseries\color{accent}Status: Solved\par}
\vspace{4pt}
{\color{accent}\hrule height 0.8pt}
\vspace{6pt}
\textbf{Area:} complexity of piecewise linear systems

\textbf{Rating rationale:} Challenging reflects a complexity
classification for recognizing an exponential finite solution count;
specialist impact concerns the solution geometry of absolute value
equations.

\subsection{Independently reviewed resolution -
2026-09-11}\label{independently-reviewed-resolution---2026-09-11}

\textbf{Affirmative complexity classification.} Theorem 1 and the
algorithm in Section 4 prove polynomial-time recognition of exactly
\(2^n\) distinct solutions to \(Ax+|x|=b\) using \(n+1\) rational LP
feasibility tests. The result uses rational binary input, has no
regularity or finiteness promise, and rejects infinite solution sets. It
classifies the displayed decision problem in \(\mathsf P\).

\textbf{Author:} Matthew J. Colbrook, Department of Applied Mathematics
and Theoretical Physics, University of Cambridge. See the
\href{https://github.com/ajt60gaibb/OpenProblemsInNLA/blob/main/references/colbrook-intervals-2026-09-11/manuscripts/AV-01.pdf}{complete
manuscript},
\href{https://github.com/ajt60gaibb/OpenProblemsInNLA/blob/main/references/colbrook-intervals-2026-09-11/verification/reviews/AV-01-review.md}{independent
agent review} and
\href{https://github.com/ajt60gaibb/OpenProblemsInNLA/blob/main/references/colbrook-intervals-2026-09-11/README.md}{submission
record}. The source archive identifies the drafts as AI-generated;
authorship is recorded at the submitter's request. This is agent
verification, not external human peer review or formal proof-assistant
certification.

The difficulty, importance and rating rationale below are historical
assessments of the original open target. The original statement and
dated audits are preserved.

\subsection{Context and notation}\label{context-and-notation}

Absolute values and vector inequalities are componentwise. All
algorithmic questions use rational input in binary and the deterministic
Turing model. Input length includes the matrix dimensions and the bit
lengths of all rational coefficients.

\subsection{Problem statement}\label{problem-statement}

For input \(A\in\mathbb Q^{n\times n}\) and \(b\in\mathbb Q^n\),
\(n\ge1\), set

\[
\Sigma_+(A,b)=\{x\in\mathbb R^n:Ax+|x|=b\}.
\]

\subsubsection{Question}\label{question}

Classify the computational complexity of deciding whether

\[
|\Sigma_+(A,b)|=2^n.
\]

In particular, is this decision problem in \(\mathsf P\), or can an
\(\mathsf{NP}\)-hardness or \(\mathsf{coNP}\)-hardness classification be
established by polynomial-time many-one reductions? The task returns one
Boolean answer. Infinite solution sets are negative instances. There is
no promise that the solution set is finite.

\newpage
\subsection{Reference}\label{reference}

Milan Hladík,
\href{https://doi.org/10.1007/s11590-025-02251-z}{\emph{Absolute value
equations with \(2^n\) solutions}}, Optimization Letters \textbf{20}
(2026), 559--575, §2.3 and §7. The paper provides structural
characterizations but explicitly leaves this recognition complexity
open. Proposition 9 treats the subclass \(\rho(|A|)<1\).

\subsection{Earlier status evidence ---
2026-09-08}\label{earlier-status-evidence-2026-09-08}

The version of record appeared October 6, 2025, in the April 2026 issue.
Searches for the title with ``complexity'', ``solved'', and ``2026'',
and for ``Hladík \(2^n\) complexity'', found no subsequent
classification. No separate arXiv version was located. The related
result about more than \(2^{n-1}\) solutions concerns a different
threshold and does not settle this question.

\subsection{Audit update --- 2026-09-10}\label{audit-update-2026-09-10}

Rechecked Hladík's
\href{https://link.springer.com/article/10.1007/s11590-025-02251-z}{journal
paper}, §2.3, which explicitly leaves maximum-count recognition open.
Searches for later complexity results found no classification of the
displayed problem. Its theorem on having more than \(2^{n-1}\) solutions
addresses a different threshold and does not solve exact \(2^n\)
recognition.
\end{document}
